% !TEX program = xelatex
% !TeX encoding = UTF-8
\documentclass[12pt,a4paper]{article}

% ============================================================
% 한국어 설정 (XeLaTeX + xeCJK)
% ============================================================
\usepackage{fontspec}
\setmainfont{Times New Roman} % 또는 Liberation Serif
\usepackage{xeCJK}
% 한국어 폰트 – Windows에 기본 설치된 Malgun Gothic 사용
\setCJKmainfont{Malgun Gothic}
\setCJKsansfont{Malgun Gothic}
\setCJKmonofont{Malgun Gothic}

% ============================================================
% 수학 및 서식
% ============================================================
\usepackage{amsmath,amssymb,amsfonts,mathtools}
\usepackage{geometry}
\geometry{left=1.25cm,right=1.25cm,top=1.25cm,bottom=3.1cm}
\usepackage{booktabs}
\usepackage{tabularx}
\usepackage{array}
\usepackage{hyperref}
\usepackage{url}
\usepackage{graphicx}
\usepackage{xcolor}
\usepackage{titlesec}
\usepackage{fancyhdr}
\usepackage{microtype}
\usepackage{multicol}
\usepackage{tikz}
\usetikzlibrary{positioning, arrows.meta, shapes.geometric, calc}

% ============================================================
% YUCT 색상 및 매크로
% ============================================================
\definecolor{skyblue}{RGB}{0,102,204}
\definecolor{darkblue}{RGB}{0,0,139}
\definecolor{gold}{RGB}{255,215,0}

\newcommand{\Keff}{K_{\mathrm{eff}}}
\newcommand{\kappac}{\kappa_c}
\newcommand{\eps}{\varepsilon}
\newcommand{\betaf}{\beta}
\newcommand{\df}{d_{\!f}}
\newcommand{\Sodd}{S_{\mathrm{odd}}}
\newcommand{\Seven}{S_{\mathrm{even}}}
\newcommand{\Mpl}{M_{\mathrm{Pl}}}
\newcommand{\Lpl}{l_{\mathrm{Pl}}}
\newcommand{\piYUCT}{\pi_{\mathrm{YUCT}}}

% ============================================================
% 절 제목 서식
% ============================================================
\titleformat{\section}{\LARGE\bfseries\color{darkblue}}{\thesection}{1em}{}
\titleformat{\subsection}{\Large\bfseries\color{darkblue}}{\thesubsection}{1em}{}

% ============================================================
% 머리말/꼬리말 (영어, 성공한 버전과 동일)
% ============================================================
\fancypagestyle{firstpage}{%
	\fancyhf{}%
	\renewcommand{\headrulewidth}{0pt}%
	\renewcommand{\footrulewidth}{0pt}%
	\fancyfoot[C]{%
		\begin{minipage}[t]{\textwidth}
			\centering
			\textcolor{gold}{\rule{\textwidth}{1.6pt}}\\[5pt]
			{\small \textsuperscript{1}Yakushev Research, \url{https://yuct.org/}} \hfill
			{\small YPSDC Protocol: \url{https://ypsdc.com/}}\\[-2pt]
			\textcolor{gold}{\rule{\textwidth}{1.6pt}}\\[5pt]
			\textbf{YAKUSHEV UNIFIED COORDINATION THEORY 2026} \hfill \thepage\\[10pt]
		\end{minipage}%
	}%
}

\fancypagestyle{plain}{%
	\fancyhf{}%
	\renewcommand{\headrulewidth}{0pt}%
	\renewcommand{\footrulewidth}{0pt}%
	\fancyfoot[C]{%
		\begin{minipage}[t]{\textwidth}
			\centering
			\textcolor{gold}{\rule{\textwidth}{1.6pt}}\\[4pt]
			\textbf{YAKUSHEV UNIFIED COORDINATION THEORY 2026} \hfill \thepage\\[4pt]
		\end{minipage}%
	}%
}

\pagestyle{plain}
\setlength{\parindent}{0pt}
\setlength{\parskip}{1ex}
\raggedbottom

\hypersetup{
	colorlinks=true,
	linkcolor=blue,
	citecolor=blue,
	urlcolor=blue,
}

% ============================================================
% 머리말 매크로 (영어, 성공한 버전과 동일)
% ============================================================
\newcommand{\makeheader}{%
	\noindent
	\begin{minipage}{0.17\textwidth}
		\raggedright
		{\fontspec{Segoe UI}[BoldFont=Segoe UI Bold]\bfseries\color{skyblue}\fontsize{46}{46}\selectfont YUCT}%
	\end{minipage}%
	\hspace{30pt}
	\begin{minipage}{0.32\textwidth}
		\raggedright
		{\sffamily\fontsize{11}{13}\selectfont YAKUSHEV\\ UNIFIED\\ COORDINATION THEORY}%
	\end{minipage}%
	\hfill
	\begin{minipage}{0.38\textwidth}
		\raggedleft
		{\sffamily\bfseries Technical Note}\\ 
		{\sffamily Version 2.3 / June 2026}\\ 
		{\sffamily\footnotesize \mbox{\url{https://doi.org/10.5281/zenodo.18444598}}}%
	\end{minipage}
	\par\vspace{5pt}
	\noindent\textcolor{gold}{\rule{\textwidth}{1.6pt}}
	\par\vspace{0pt}
}

% ============================================================
% 문서 시작
% ============================================================
\begin{document}
	
	% 표지
	\thispagestyle{firstpage}
	\makeheader
	
	\vspace{0cm}
	{\LARGE\bfseries 부록 AF: YUCT에서의 현실의 대수적 프레임워크 \\ 
		\Large 질량, 기하, 시간, 생명, 정보, 소수, 사회를 지배하는 22개의 닫힌 루프로 이루어진 자기일관적 체계 \par}
	\vspace{0.2cm}
	
	{\large 알렉세이 V. 야쿠셰프\textsuperscript{1}} \\
	\vspace{0.0cm}
	
	{\color{darkblue}\textbf{\LARGE 초록}} \par
	{\large
		\noindent
		야쿠셰프 통일 조정 이론(YUCT)은 독립적인 경험적 피팅의 집합이 아니라, 상호 연결된 \textbf{대수적 루프}들의 집합으로 정의되는 엄밀한 수학적 구조이다. 본 부록에서는 이론의 공리적 골격을 형성하는 22개의 주요 루프를 집약한다. 우리는 조정의 계층적 자기유사성으로부터 도출된 보편 상수 \(S_{\mathrm{odd}}=6/5\)와 \(S_{\mathrm{even}}=4/5\)가 임의의 입력이 아니라 닫힌 자기지시적 시스템의 필연적 결과임을 보여준다. 이 루프들은 소립자의 질량과 시공간의 기하로부터 유전체의 구조와 동역학, 경제 성장, 양자 얽힘의 원리, 소수의 분포, 안전한 정보 프로토콜의 설계에 이르기까지 모든 규모의 현상을 통합한다. 우리는 \textbf{하드 루프}(보편 불변량)와 \textbf{소프트 루프}(특정 시스템에 적용되는 보편 법칙) 사이의 결정적 구분을 도입하여 겉보기 편차를 해결하고 프레임워크의 공학적 역량을 강조한다. 이 다중 루프 대수적 프레임워크의 존재——자유 연속 매개변수가 없음——는 YUCT를 현상론적 모델에서 물리적, 생물학적, 수학적, 정보적 현실의 반증 가능한 통일 골격으로 격상시킨다.}
	
	\vspace{0.1cm}
	\noindent
	{\color{darkblue}\textbf{키워드:}} YUCT, 대수적 루프, 조정 효율, 보편 상수, 페르미온 질량 계층, 창발적 기하, 시간 양자화, 게놈 구조, 대사 스케일링, 양자 얽힘, 정보 이론, 소수, 자기일관성, 반증 가능성.
	\vspace{0.1cm}
	\noindent
	{\color{darkblue}\textbf{인용:}} Yakushev AV. 부록 AF: YUCT에서의 현실의 대수적 프레임워크. 2026. DOI: \url{https://doi.org/10.5281/zenodo.18444598} (본 권).
	
	\vspace{0.4cm}
	
	\begin{multicols}{2}
		
		\section{기초: 조정, D+I·R, 보편 스케일링 법칙}
		
		야쿠셰프 통일 조정 이론(YUCT)은 단일 존재론적 원시 개념에 기반한다: \textbf{조정}——분산 시스템의 구성 요소들이 그 상태를 정렬하는 과정. 기본 메커니즘은 YPSDC 프로토콜(야쿠셰프 동기 분산 조정 프로토콜)이며, 이는 오프라인 단계(공유 \textit{사전} \(\mathcal{D}\)의 배포)를 온라인 단계(짧은 \textit{색인} \(\kappa\)를 전송하여 사전 합의된 복잡한 동작을 활성화)로부터 분리한다.
		
		임의의 조정 시스템의 상태는 \textbf{D+I·R 삼항}에 의해 기술된다:
		\[
		\text{현실} = \mathbf{D} + \mathbf{I} \times \mathbf{R},
		\]
		여기서 \(\mathbf{D}\)는 사전(가능한 상태와 규칙의 집합), \(\mathbf{I}\)는 정보(실현된 상태, 엔트로피로 측정), \(\mathbf{R} \ge 1\)은 공명(코히런스 증폭 인자)이다.
		
		\textbf{조정 효율}은 정보 이론적 압축비로 정의된다:
		\[
		\Keff = \frac{H(\mathcal{A})}{H(\kappa)} = \frac{\text{행동 엔트로피}}{\text{색인 엔트로피}} \ge 1.
		\]
		YUCT의 기본 결과는 \textbf{보편 스케일링 법칙}이며, 이는 임의의 조정 과정에서의 상대 오차(또는 변동) \(\varepsilon\)을 기술한다:
		\[
		\boxed{\varepsilon = \kappac \, \alpha \, \Keff^{-\beta}, \qquad \beta = \frac{2}{3}, \quad \kappac = \frac{1}{3}},
		\]
		여기서 \(\alpha\)는 시스템 의존 상수이다. 이 법칙은 화학 결합 에너지에서 우주 마이크로파 배경 변동, 그리고 최근 발견된 소수 분포(부록~PrimeN)에 이르기까지 50자리 이상에 걸쳐 경험적으로 검증되었다. 지수 \(\beta = 2/3\)는 경험적 피팅이 아니라 아래에서 유도되는 닫힌 대수 구조의 필연적 결과이다.
		
		\section{일차 대수적 루프: \(S_{\mathrm{odd}}, S_{\mathrm{even}}\)와 \(\beta\)의 유도}
		
		조정 네트워크의 계층적 자기유사성은 연속적인 수준의 기여가 보편 지수 \(\beta\)로 기하학적으로 스케일링됨을 의미한다. 무한 계층 위의 기하급수를 합하면, 단방향(홀수) 및 교대(짝수) 상관의 극한 기여를 얻는다:
		\begin{equation}
			S_{\mathrm{odd}} = \sum_{k=0}^{\infty} \beta^{2k+1} = \frac{\beta}{1-\beta^{2}}, \qquad
			S_{\mathrm{even}} = \sum_{k=1}^{\infty} \beta^{2k} = \frac{\beta^{2}}{1-\beta^{2}}.
		\end{equation}
		이 정의들은 즉시 두 개의 정확한 관계를 제공한다:
		\begin{equation}
			S_{\mathrm{odd}} + S_{\mathrm{even}} = 2, \qquad \frac{S_{\mathrm{odd}}}{S_{\mathrm{even}}} = \frac{1}{\beta}.
		\end{equation}
		시스템이 내부적으로 자기일관적이고 관측되는 물리적 조정 네트워크의 프랙탈 차원 (\(d_f \approx 2.2\))에 대응하려면, 지수는 다음 값을 취해야 한다:
		\begin{equation}
			\boxed{\beta = \frac{2}{3}}.
		\end{equation}
		\(\beta = 2/3\)을 루프 방정식에 대입하면, 이후의 모든 유도의 기초가 되는 정확한 유리 상수를 얻는다:
		\begin{equation}
			\boxed{S_{\mathrm{odd}} = \frac{6}{5} = 1.2, \qquad S_{\mathrm{even}} = \frac{4}{5} = 0.8.}
		\end{equation}
		이것이 \textbf{일차 대수적 루프}이다: 세 개의 수 (\(\beta, S_{\mathrm{odd}}, S_{\mathrm{even}}\))가 자유 매개변수 없이 불가분하게 결합되어 있다. 비 \(S_{\mathrm{odd}}/S_{\mathrm{even}} = 3/2\)는 기본적인 세대 간 스케일링 인자이다.
		
		\section{루프 II: 페르미온 질량의 반정수 양자화}
		
		일차 루프는 세대 간 질량비가 인자 \(3/2 = S_{\mathrm{odd}}/S_{\mathrm{even}}\)에 의해 지배됨을 결정한다. 지수 \(\beta = 2/3\) 자체는 \(2 \times (1/3)\)로 분해되며, \(1/3\)은 세 개의 공간 차원을, \(2\)는 스핀 통계를 반영한다. 이는 진정한 \textbf{스케일링 양자}를 드러낸다:
		\begin{equation}
			q \equiv (3/2)^{1/3} = \left( \frac{S_{\mathrm{odd}}}{S_{\mathrm{even}}} \right)^{1/3} \approx 1.1447.
		\end{equation}
		아홉 종류의 하전 페르미온의 질량은 모두 \(\ln q\)의 반정수 단위로 양자화된다:
		\begin{equation}
			m_f = m_e \cdot q^{N_f} \cdot \eta_f, \qquad N_f \in \frac{1}{2}\mathbb{Z}.
		\end{equation}
		실험 데이터와의 서브퍼센트 일치는 일차 대수적 루프의 폐쇄의 직접적 결과이다.
		
		\section{루프 III: 약한 스케일 계층}
		
		약한 스케일은 조정 깊이 \(L = 96\)에 의해 생성되며, 이는 표준 모형에서의 바일 페르미온 자유도의 총수(우측 중성미자 포함)에 의해 고정된다: \(96 = 3 \times 16 \times 2\). 표준 플랑크 질량 \(M_{\mathrm{Pl}} \approx 1.2209 \times 10^{19}\)~GeV, 결합 인자 \(\kappa_c = 1/3\), 그리고 YUCT에서 창발되는 \(\pi\)의 값 \(\pi_{\mathrm{YUCT}} = S_{\mathrm{odd}}\varphi^2\)을 사용하면, \(W\) 보손의 질량은:
		\begin{equation}
			\boxed{m_W = \left( \kappa_c \, \pi_{\mathrm{YUCT}} \right)^{-1/2} \; M_{\mathrm{Pl}} \; \left( \frac{2}{3} \right)^{96}},
		\end{equation}
		여기서 \(\pi_{\mathrm{YUCT}} = S_{\mathrm{odd}}\varphi^2 \approx 3.1416407865\). 수치적으로는:
	\end{multicols}
	\small
	\[
	\left( \frac{1}{3} \times 3.1416408 \right)^{-1/2} \approx (1.0472136)^{-1/2} \approx 0.977,\qquad 
	(2/3)^{96} \approx 6.75 \times 10^{-18},
	\]
	\[
	m_W \approx 0.977 \times 1.2209 \times 10^{19} \times 6.75 \times 10^{-18} \approx 80.5\ \text{GeV}.
	\]
	이 루프는 일차 루프의 상수와 이산 정수 \(96\)을 통해 플랑크 스케일을 전약한 스케일에 연결한다. 미세 조정은 필요 없으며, 루프 V와 동일한 플랑크 질량이 사용된다.
	\begin{multicols}{2}
		\section{루프 IV: 기하의 창발과 \(\pi\)–\(\varphi\) 연결}
		
		황금비 \(\varphi = (1+\sqrt{5})/2\)(프랙탈 계층의 기하학적 불변량)을 사용하면:
		\begin{equation}
			S_{\mathrm{odd}} \cdot \varphi^2 = \frac{6}{5} \left( \frac{1+\sqrt{5}}{2} \right)^2 = \frac{9+3\sqrt{5}}{5} \approx 3.1416407865.
		\end{equation}
		이 값은 \(\pi \approx 3.1415926535\)를 상대 오차 \(1.5 \times 10^{-5}\)로 근사한다. YUCT에서 \(\pi\)는 \(\Keff \to \infty\)의 점근 극한이며, 유한 시스템에서는 원주와 지름의 비가 \(S_{\mathrm{odd}}\varphi^2\)로 재규격화된다.
		
		\section{루프 V: 우주의 중입자 질량}
		
		관측 가능 우주의 총 중입자 질량 \(M_{\mathrm{bar}}\)과 양성자 질량 \(m_p\)의 비는:
		\begin{equation}
			\frac{M_{\mathrm{bar}}}{m_p} = \left( \frac{M_{\mathrm{Pl}}}{m_e} \right)^{\mathcal{S}},
		\end{equation}
		여기서 지수 \(\mathcal{S}\)는 기본 대수 상수의 합이다:
		\begin{equation}
			\mathcal{S} \equiv S_{\mathrm{odd}} + S_{\mathrm{even}} + \frac{S_{\mathrm{odd}}}{S_{\mathrm{even}}} = \frac{6}{5} + \frac{4}{5} + \frac{3}{2} = 3.5.
		\end{equation}
		여기서 \(M_{\mathrm{Pl}}\)은 루프 III과 동일한 표준 플랑크 질량이다. 수치적으로 \((M_{\mathrm{Pl}}/m_e)^{3.5} \approx 1.88 \times 10^{78}\)이며, 관측값 \(M_{\mathrm{bar}}/m_p \approx 1.4 \times 10^{78}\)과 인자 1.3 이내로 일치한다.
		
		\section{루프 VI: 시간의 양자화}
		
		조정 엔트로피는 \(S_0 = k_B \ln 2\)의 단위로 양자화된다. 따라서 고유 시간은 이산적 단계로 진행된다. 최소 시간 양자와 플랑크 시간의 비는 동일한 계층 상수에 의해 고정된다:
		\begin{equation}
			\boxed{\frac{\Delta \tau_{\min}}{t_P} = \frac{S_{\mathrm{even}}}{S_{\mathrm{odd}}} = \beta = \frac{2}{3}}.
		\end{equation}
		거시적으로, 질량 \(M\)의 블랙홀의 시간의 상대 변동은 다음과 같이 스케일된다:
		\begin{equation}
			\frac{\delta \tau}{\tau} \propto M^{-\beta} = M^{-0.67}.
		\end{equation}
		이 예측은 준주기 진동(QPO) 및 중력파 링잉 감쇠(부록~G, 부록~V 참조)를 통해 검증 가능하다.
		
		\section{루프 XXI: 파이겐바움 상수와 질서–혼돈의 다리}
		
		혼돈 이론의 보편 상수——파이겐바움 상수
		\(\delta \approx 4.6692016\)(주기 배가 분기 매개변수)와
		\(\alpha \approx 2.502907875\)(스케일 매개변수)——는 독립적인 경험 수가 아니라 YUCT의 대수적 골격과 강성적으로 연결되어 있다.
		이산적 조정 질서(지수 \(\beta=2/3\))와 연속적 재규격화 군 캐스케이드(혼돈) 사이의 다리는 정확한 해석적 관계
		\begin{equation}
			2\alpha^2 = \pi \delta,
			\label{eq:feigenbaum-pi-exact}
		\end{equation}
		에 의해 주어지며, 이는 복소 평면 상의 주기 배가 작용소 아래의 재규격화 군 구조의 엄밀한 결과이다~\cite{Feigenbaum1978}.
		(\ref{eq:feigenbaum-pi-exact})와 YUCT에서 유도된 \(\pi\)의 고정밀 근사(루프~IV)를 결합하면,
		\begin{equation}
			\pi \;\approx\; \pi_{\mathrm{YUCT}} = S_{\mathrm{odd}}\,\varphi^2,
			\qquad
			S_{\mathrm{odd}} = \frac{6}{5},\;\; \varphi = \frac{1+\sqrt{5}}{2},
		\end{equation}
		닫힌 표현을 얻는다:
		\begin{equation}
			\boxed{\;
				2\alpha^2 \;\approx\; \bigl(S_{\mathrm{odd}}\,\varphi^2\bigr)\,\delta
				\;}.
			\label{eq:feigenbaum-yuct}
		\end{equation}
		방정식~(\ref{eq:feigenbaum-yuct})는 질서의 보편 상수 (\(S_{\mathrm{odd}}, S_{\mathrm{even}}, \beta\))와 혼돈의 보편 상수 (\(\alpha, \delta\)) 사이에 직접적인 매개변수-프리 대수적 연결을 확립한다. 작은 수치적 편차
		\(\Delta\pi = |\pi - \pi_{\mathrm{YUCT}}| \approx 4.8\times10^{-5}\)
		는 유도의 결함이 아니라, 물리적 기하의 유한한 조정 효율을 반영한 YUCT의 \textbf{반증 가능한 예측}이다(부록~\ref{sec:unity-of-reality} 및~\ref{sec:ehrenfest} 참조).
		
		이 루프는 창조(질서)의 상수와 파괴(혼돈)의 상수가 독립적이 아니라 단일한 통일된 수학적 현실의 두 측면임을 보여준다. 질서–혼돈의 다리는 세포 자동자 규칙~30의 분석(부록~UoR)에 의해 더욱 강화되며, 결정론적 혼돈이 완벽하게 혼합되는 임계 깊이가 관계 \(\ln t_{\mathrm{crit}} \sim \delta/\varphi\)를 통해 파이겐바움 상수와 대수적으로 관련됨이 보여진다. YUCT 라그랑지안으로부터 파이겐바움 상수의 완전한 유도는 미해결 과제이지만, 여기서 제시된 구조적 연결은 수치적 일치를 이론의 검증 가능한 구성 요소로 변환한다.
		
		\section{루프 XXII: 소수 조정 사다리}
		
		상수 \(S_{\mathrm{odd}}\)와 \(S_{\mathrm{even}}\)는 소수(오랫동안 순수 수학으로 간주되어 온 영역)의 변동도 지배한다. \(n\)번째 소수 \(p_n\)를 \(K_{\mathrm{eff}} \propto p_n\)인 조정 사다리 위의 노드로 해석하면, 보편 오차 법칙은 고전적 Rosser 근사로부터의 상대 편차가 \(p_n^{-2/3}\)로 스케일됨을 예측한다. \(n = 5\times10^5\)까지의 수치 해석은 국소 스케일링 지수가 \(2/3\)로 수렴함을 보여준다(점근값 \(0.66666\)).
		
		또한, 대수적 루프는 Rosser 공식에 대한 \textbf{매개변수-프리 보정}을 제공한다:
		\[
		p_n \approx R_n - \frac{S_{\mathrm{even}}}{2}\, n^{1-\beta}\,\ln n,
		\]
		이로 인해 최대 오차가 \(\approx 2.3\%\)(보통 Rosser 공식)에서 \(n\le 10^5\)에 대해 \(\approx 0.012\%\)로 감소한다. 개선된 경험적 보정(\(A = -0.44\), \(B = 1.05\))은 이를 \(\approx 0.006\%\)로 개선한다.
		
		이 루프는 경험적 보정 상수 \(A\)와 \(B\)가 아직 제1원리로부터 유도되지 않았기 때문에 \textbf{소프트}로 분류된다. 그러나 이론적 형식은 이미 자유 매개변수를 포함하지 않으며 대수적 골격의 보편성을 입증한다. 완전한 유도는 부록~PrimeN을 참조.
		
	\end{multicols}
	
	\section*{YUCT의 여섯 가지 코어 물리 루프}
	
	\begin{table}[htbp]
		\centering
		\caption{완전한 대수적 프레임워크: 자유 매개변수 없는 여섯 개의 상호 연결된 코어 루프.}
		\label{tab:all-loops}
		\footnotesize
		\begin{tabular}{l c c}
			\toprule
			\textbf{루프} & \textbf{코어 관계} & \textbf{검증 근거} \\
			\midrule
			I. 일차 상수 & \(S_{\mathrm{odd}}=1.2, S_{\mathrm{even}}=0.8, \beta=2/3\) & 프랙탈 기하, KPZ 관계 \\
			II. 페르미온 질량 & \(m_f = m_e q^{N_f}, q=(3/2)^{1/3}\) & PDG 데이터 (오차 <1\%) \\
			III. 약한 스케일 & \(m_W = (\kappa_c \pi_{\mathrm{YUCT}})^{-1/2} M_{\mathrm{Pl}} (2/3)^{96}\) & \(m_W = 80.4\) GeV \\
			IV. 창발적 기하 & \(S_{\mathrm{odd}}\varphi^2 \approx \pi\) (오차 \(10^{-5}\)) & 수학적 자기일관성 \\
			V. 우주 질량 & \(M_{\mathrm{bar}}/m_p = (M_{\mathrm{Pl}}/m_e)^{3.5}\) & Planck 2018, \(H_0\) 측정 \\
			VI. 시간 양자화 & \(\Delta \tau_{\min} = \frac{2}{3} t_P, \ \frac{\delta \tau}{\tau} \propto M^{-0.67}\) & GWTC-4.0 링잉 스케일링 \\
			\bottomrule
		\end{tabular}
	\end{table}
	
	\section*{확장 인벤토리: 과학에서의 22개 대수적 루프}
	
	위의 여섯 개의 코어 루프는 물리적 현실의 골격을 형성한다. 그러나 동일한 대수 구조——상수 \(S_{\mathrm{odd}} = 1.2\), \(S_{\mathrm{even}} = 0.8\), \(\beta = 2/3\)에 뿌리를 둔——는 광범위한 과학 기술 영역에 투영된다. 표~\ref{tab:all-loops-22}는 YUCT 부록에서 확인된 22개의 대수적 루프의 완전한 인벤토리이며, 물리, 화학, 생물, 경제, 사회, 정보 이론, 양자 역학, 수론을 망라한다. 각 루프는 일차 대수적 폐쇄의 직접적 결과이며, 프레임워크에 독립적인 경험적 검증을 제공한다.
	
	% 조정된 표 – \footnotesize, 좁은 열 간격, 간결한 텍스트
	\begin{table}[htbp]
		\centering
		\footnotesize
		\setlength{\tabcolsep}{2.5pt}
		\caption{YUCT에서의 22개 대수적 루프의 완전 인벤토리 (XXII: 소수로 업데이트)}
		\label{tab:all-loops-22}
		\begin{tabular}{l c c p{3.8cm} c}
			\toprule
			\textbf{루프} & \textbf{클래스} & \textbf{영역} & \textbf{관계식} & \textbf{부록} \\
			\midrule
			I   & 하드 & 기본 상수 & \(S_{\mathrm{odd}}=1.2,\ S_{\mathrm{even}}=0.8,\ \beta=2/3\) & Ferra1.2 \\
			II  & 하드 & 입자 질량 & \(m_f = m_e q^{N_f},\ q=(3/2)^{1/3}\) & J, \(\Lambda\) \\
			III & 하드 & 약한 스케일 & \(m_W = (\kappa_c \pi_{\mathrm{YUCT}})^{-1/2} M_{\mathrm{Pl}} (2/3)^{96}\) & \(\Lambda\) \\
			IV  & 하드 & 기하 & \(S_{\mathrm{odd}}\varphi^2 \approx \pi\) (오차 \(10^{-5}\)) & Ferra1.2, Ehrenfest \\
			V   & 하드 & 우주 질량 & \(M_{\mathrm{bar}}/m_p = (M_{\mathrm{Pl}}/m_e)^{3.5}\) & \(\Omega\), \(\Lambda\) \\
			VI  & 하드 & 시간 양자화 & \(\Delta\tau_{\min} = \frac{2}{3}t_P,\ \frac{\delta\tau}{\tau} \propto M^{-0.67}\) & V, G \\
			VII & 소프트 & 화학 결합 & \(E_{\mathrm{bond}} \propto r^{-2/3}\) & C, Z \\
			VIII& 소프트 & 대사 & \(B \propto M^{\beta d_f/2}\) & D, E.7 \\
			IX  & 소프트 & 변이율 & \(\mu \propto K_{\mathrm{repair}}^{-2/3}\) & E \\
			X   & 소프트 & 경제 & \(\sigma_{\mathrm{GDP}} \propto W^{-2/3}\) & F \\
			XI  & 소프트 & 사회 집단 & \(N_{\mathrm{crit}}/N_{\mathrm{opt}} \approx 1.5\) & O \\
			XII & 소프트 & 진화적 특이점 & 쌍곡적 성장, \(t_s \approx 2045\) & T \\
			XIII& 소프트 & 상온 핵융합 & \(K_{\mathrm{crit}} \sim 10^{12}\), 코히/인코히 = 1.5 & R \\
			XIV & 하드 & 게놈 구조 & \(\frac{\text{AA+TT+GG+CC}}{\text{AT+TA+GC+CG}} \approx 1.5\) & E \\
			XV  & 소프트 & 크로마틴 & \(S \propto L^{0.733}\) (\(d_f \approx 2.2\)) & E, D \\
			XVI & 하드 & 코돈 사용 & 빈도 어트랙터 \(\varphi \approx 1.618\) & E \\
			XVII& 소프트 & 인구 임계 & \(N_{\mathrm{crit}}/N_{\mathrm{opt}} \approx 1.5\) (던바) & O, E \\
			XVIII&소프트 & 진화적 템포 & 종분화율 \(\propto K_{\mathrm{eff}}^{-0.67}\) & T, E \\
			XIX & 소프트 & 양자 얽힘 & \(K_{\mathrm{eff}} \to \infty,\ S_{\mathrm{QM}} = 2\sqrt{2}\) & G, X, Y \\
			XX  & 소프트 & 2요소 인증 & \(K_{\mathrm{eff}} = 160/20 = 8\) & B, K, X \\
			XXI & 소프트 & 혼돈 이론 & \(\delta = \frac{S_{\mathrm{odd}}}{S_{\mathrm{even}}} \pi_{\mathrm{YUCT}} \ln(\alpha/\beta)\) & 현실의 통일 \\
			\textbf{XXII} & \textbf{소프트} & \textbf{소수} & \(\mathbf{p_n \approx R_n - \frac{S_{\mathrm{even}}}{2}n^{1-\beta}\ln n}\) & \textbf{PrimeN} \\
			\bottomrule
		\end{tabular}
	\end{table}
	
	\begin{figure}[htbp]
		\centering
		\begin{tikzpicture}[scale=0.9, transform shape]
			\tikzset{
				loop/.style={circle, draw=blue!70, fill=blue!20, minimum size=0.8cm, inner sep=0pt, font=\scriptsize},
				label/.style={font=\scriptsize, align=center},
				axis/.style={->, thick, gray},
				domainline/.style={thin, dashed, gray!70},
			}
			
			% 축
			\draw[axis] (0,0) -- (16,0) node[right] {\scriptsize 영역 스케일 (자릿수)};
			\draw[axis] (0,-2.5) -- (0,4.5) node[above] {\scriptsize 정밀도 / 일치도};
			
			% Y축 눈금
			\foreach \y/\ylab in {-1.5/10^{-15}, 0/10^0, 1.5/10^{15}, 3/10^{30}} {
				\draw (-0.1,\y) -- (0.1,\y) node[left] {\scriptsize $\ylab$};
			}
			
			% X축 눈금
			\foreach \x/\xlab in {2/원자, 5/생물, 8/사회, 11/우주} {
				\draw (\x,0.1) -- (\x,-0.1) node[below] {\scriptsize \xlab};
			}
			
			% 배경 영역선
			\draw[domainline] (2,-2.5) -- (2,4);
			\draw[domainline] (5,-2.5) -- (5,4);
			\draw[domainline] (8,-2.5) -- (8,4);
			\draw[domainline] (11,-2.5) -- (11,4);
			
			% 루프 좌표 (x=스케일, y=정밀도)
			% 코어 물리 (I-VI)
			\node[loop] (L1) at (0.5, 3.8) {I};
			\node[label, above=0.2cm of L1] {일차};
			\node[loop] (L2) at (3.0, 3.4) {II};
			\node[label, above=0.2cm of L2] {페르미온};
			\node[loop] (L3) at (5.5, 3.0) {III};
			\node[label, above=0.2cm of L3] {약함};
			\node[loop] (L4) at (7.0, 2.6) {IV};
			\node[label, above=0.2cm of L4] {기하};
			\node[loop] (L5) at (10.0, 2.2) {V};
			\node[label, above=0.2cm of L5] {중입자};
			\node[loop] (L6) at (12.0, 1.8) {VI};
			\node[label, above=0.2cm of L6] {시간};
			
			% 학제 (VII-XIII)
			\node[loop] (L7) at (2.0, 1.2) {VII};
			\node[label, below=0.2cm of L7] {화학결합};
			\node[loop] (L8) at (4.5, 0.8) {VIII};
			\node[label, below=0.2cm of L8] {대사};
			\node[loop] (L9) at (6.5, 0.4) {IX};
			\node[label, below=0.2cm of L9] {돌연변이};
			\node[loop] (L10) at (8.5, -0.2) {X};
			\node[label, below=0.2cm of L10] {경제};
			\node[loop] (L11) at (3.5, -0.8) {XI};
			\node[label, below=0.2cm of L11] {사회};
			\node[loop] (L12) at (5.5, -1.2) {XII};
			\node[label, below=0.2cm of L12] {진화};
			\node[loop] (L13) at (7.5, -1.6) {XIII};
			\node[label, below=0.2cm of L13] {상온핵융합};
			
			% 유전자·생물 (XIV-XVIII)
			\node[loop] (L14) at (1.5, 2.2) {XIV};
			\node[label, right=0.2cm of L14] {게놈};
			\node[loop] (L15) at (3.5, 2.0) {XV};
			\node[label, right=0.2cm of L15] {크로마틴};
			\node[loop] (L16) at (5.5, 1.8) {XVI};
			\node[label, right=0.2cm of L16] {코돈};
			\node[loop] (L17) at (2.5, -0.2) {XVII};
			\node[label, right=0.2cm of L17] {인구};
			\node[loop] (L18) at (6.5, -0.8) {XVIII};
			\node[label, right=0.2cm of L18] {종분화};
			
			% 정보·양자 (XIX-XX)
			\node[loop, draw=red!70, fill=red!20] (L19) at (7.0, 2.5) {XIX};
			\node[label, above=-1.1cm of L19] {EPR 얽힘};
			\node[loop, draw=red!70, fill=red!20] (L20) at (8.5, 0.5) {XX};
			\node[label, below=-1.2cm of L20] {2FA};
			
			% 루프 XXI (파이겐바움)
			\node[loop, draw=purple!70, fill=purple!20] (L21) at (11.5, 0.2) {XXI};
			\node[label, below=0.2cm of L21] {파이겐바움};
			
			% 신규: 루프 XXII (소수)
			\node[loop, draw=orange!80, fill=orange!20] (L22) at (10.0, -1.6) {XXII};
			\node[label, below=0.2cm of L22] {소수};
			
			% 연결선
			\draw[blue!40, thick, dashed] (L1) -- (L2) -- (L3) -- (L4) -- (L5) -- (L6);
			\draw[blue!40, thick, dashed] (L1) -- (L7) -- (L8) -- (L9) -- (L10);
			\draw[blue!40, thick, dashed] (L7) -- (L11) -- (L12) -- (L13);
			\draw[green!60!black, thick, dashed] (L1) -- (L14) -- (L15) -- (L16);
			\draw[green!60!black, thick, dashed] (L14) -- (L17) -- (L18);
			\draw[red!60, thick, dashed] (L1) -- (L19);
			\draw[red!60, thick, dashed] (L19) -- (L20);
			\draw[purple!60, thick, dashed] (L4) -- (L21);
			\draw[orange!80, thick, dashed] (L1) -- (L22);
			
			% 제목
			\node[font=\bfseries\small, align=center] at (8, 5.0) {YUCT의 22개 대수적 루프\\50자리 이상에 걸친 수렴};
			
			% 범례
			\node[draw, fill=white, rounded corners, font=\tiny, align=left] at (13, 3.5) {
				\textbf{코어 물리 (I--VI)}\\
				\textbf{학제 (VII--XIII)}\\
				\textbf{유전자·생물 (XIV--XVIII)}\\
				\textbf{정보·양자 (XIX--XX)}\\
				\textbf{혼돈 이론 (XXI)}\\
				\textbf{소수 (XXII)}\\
				\textit{모든 루프는}\\ 
				\textit{$S_{\mathrm{odd}}=1.2$, $S_{\mathrm{even}}=0.8$에서 유래}
			};
		\end{tikzpicture}
		\caption{YUCT의 22개 대수적 루프의 도식. 원자에서 우주 스케일까지, 또한 학제를 넘어 확장되며, 파이겐바움 혼돈 루프(XXI)와 새로 추가된 소수 조정 사다리(XXII)를 포함. 모든 루프는 동일한 기본 상수 $S_{\mathrm{odd}} = 1.2$와 $S_{\mathrm{even}} = 0.8$에 뿌리를 두고 있다.}
		\label{fig:twentyone-loops}
	\end{figure}
	
	\vspace{0.5cm}
	\noindent
	\textbf{스펙트럼 클러스터링의 해석: 계층적 조정 층.}
	그림~\ref{fig:twentyone-loops}의 눈에 띄는 시각적 특징은 대수적 루프가 거의 수직인 ``능선''으로 클러스터링된 것으로, 이는 특정 조직 수준——원자, 생물, 사회, 우주——에 대응한다. 이는 플롯의 인공물이 아니라 YUCT의 핵심인 \textbf{계층적이고 층상의 조정 구조}의 직접적 표현이다.
	
	YUCT 프레임워크에서 현실은 연속 스펙트럼이 아니라, 각 층이 지배적 사전과 특징적 스케일에 의해 특징지어지는 이산적 조정 층의 스택이다. 기본 대수적 루프(루프 I)는 스케일 불변이지만, 특정 영역으로의 \textit{사영}은 그 층의 특징적 변수(예: 결합 에너지, 돌연변이율, 집단 크기, 우주론적 밀도)에 의존한다. 관측되는 네 개의 능선은 바로 이러한 지배적 층에 대응한다:
	\begin{itemize}
		\item \textbf{원자 스케일 (X $\approx$ 2):} 양자장과 분자 궤도에 의한 조정. 여기서의 루프는 화학 결합과 핵 과정을 지배한다.
		\item \textbf{생물 스케일 (X $\approx$ 5):} 유전 암호와 대사 네트워크에 의한 조정. 여기서의 루프는 게놈 구조, 대사, 돌연변이율을 지배한다.
		\item \textbf{사회 스케일 (X $\approx$ 8):} 언어, 문화, 경제 교환에 의한 조정. 여기서의 루프는 집단 역학, 경제 변동, 진화적 특이점을 지배한다. 이 능선은 또한 \textbf{인간이 설계한 정보 시스템}(루프 XX)을 수용한다.
		\item \textbf{우주 스케일 (X $\approx$ 11):} 중력과 전지구적 회전에 의한 조정. 여기서의 루프는 우주의 중입자 질량, 약한 스케일 계층, 시간의 양자화를 지배한다.
	\end{itemize}
	
	결정적으로, 이 구조는 \textbf{예측적}이다. YUCT가 옳다면, 이 능선들의 \textit{사이}에 있는 영역(예: 생물학과 사회의 경계의 신경과학, 또는 사회와 우주 사이의 행성 역학)에서 새로 발견되는 대수적 루프는 자연스럽게 도표 상의 해당 중간 영역에 떨어지며, 동일한 기본 상수 \(S_{\mathrm{odd}}\)와 \(S_{\mathrm{even}}\)를 유지할 것이다. 반대로, 이 클러스터화 구조의 \textit{바깥}에 있는 강건한 루프의 발견은 이론에 대한 중대한 도전이 될 것이다. 22개 루프의 스펙트럼 클러스터링은 따라서 YUCT의 계층적 존재론의 검증인 동시에 미래의 학제적 발견을 위한 로드맵이기도 하다.
	
	\vspace{0.5cm}
	\noindent
	\textbf{두 가지 클래스의 대수적 루프: 하드 제약 대 소프트 발현.}
	22개 대수적 루프의 인벤토리는 정확한 유리 항등식에서 문맥 의존적 스케일링 관계까지 이질적으로 보일 수 있다. 이 일견하는 다양성은 약점이 아니라 YUCT의 계층 구조를 반영한다. 모든 루프는 두 가지 기본 클래스 중 하나에 속하며, 그 매개변수가 하위 대수적 골격에 의해 고정되는 정도에 따라 구별된다.
	
	\paragraph{클래스 I: 하드 루프 (보편 불변량).}
	이 루프들은 \textbf{정확한 수학적 항등식}을 표현하거나 일차 대수적 루프 (\(S_{\mathrm{odd}}=1.2, S_{\mathrm{even}}=0.8, \beta=2/3\))에 의해 강성적으로 결정되는 \textbf{기본 상수}를 포함한다. 그 값들은 협상의 여지가 없으며, 어떤 경험적 편차라도 YUCT의 핵심을 즉시 반증할 것이다. 하드 루프는 ``현실의 골격''——우주의 불변 수치 문법——을 형성한다.
	\begin{itemize}
		\item \textbf{예:} 루프 I--VI, XIV, XVI.
		\item \textbf{고정량:} \(\beta = 2/3\), \(S_{\mathrm{odd}}=1.2\), \(S_{\mathrm{even}}=0.8\), \(m_W \approx 80.4\ \text{GeV}\), 근사 \(S_{\mathrm{odd}}\varphi^2 \approx \pi\), 이염기비 \(\approx 1.5\).
		\item \textbf{왜 하드인가:} 이 숫자들은 특정 관측 시스템으로부터 독립적이다. 그것들은 현실 자체의 ``운영 체제''에 내장되어 있다.
	\end{itemize}
	
	\paragraph{클래스 II: 소프트 루프 (특정 시스템에 적용되는 보편 법칙).}
	이 루프들은 \textbf{스케일링 관계} 또는 \textbf{조정 원리}를 기술하며, 그 함수 형태는 보편 상수(예: \(\varepsilon \propto K_{\mathrm{eff}}^{-2/3}\))에 의해 주어지지만, 구체적 수치(예: 특정 실험에서의 정확한 효율 \(K_{\mathrm{eff}}\))는 시스템의 우연적 성질에 의존한다. 그것들은 특정 수로부터의 편차에 의해 반증되지 않고, 예측된 스케일링 법칙의 체계적 실패 또는 예측된 \textit{최적} 상태를 달성하지 못함으로써 반증된다.
	\begin{itemize}
		\item \textbf{예:} 루프 VII--XIII, XV, XVII--XXII.
		\item \textbf{가변량:} 2FA의 정확한 \(K_{\mathrm{eff}}\)(키 길이에 의존), 화학 반응의 정확한 속도(온도에 의존), 사회 집단의 실제 크기(문화적 요인에 의존), 소수의 경험적 보정 상수.
		\item \textbf{그래도 루프인 이유:} 그것들은 \textbf{동일한} 함수 법칙(\(\propto x^{-2/3}\), 최적비 1.5)이 극히 다른 영역에서 조정을 지배함을 보여준다. 그것들은 현실의 보편적 ``동사''이며, 각 ``명사''(시스템)에 의해 다르게 활용된다.
	\end{itemize}
	
	\paragraph{편차의 조화: 수정된 조정의 경우.}
	이 분류는 사회학이나 상온 핵융합과 같은 분야에서의 겉보기 편차가 왜 루프를 파괴하지 않는지를 설명한다. 소셜 네트워크나 실험실 LENR 실험과 같은 현상은 \textbf{수정된 조정}의 예이다. 인간의 기술은 자연적 어트랙터를 일시적으로 오버라이드하거나(예: 던바 수를 초과하는 집단 생성), 외부 에너지와 구조를 주입함으로써 인공적으로 상태를 달성하려고 시도할 수 있다(예: 상온 핵융합의 \(K_{\mathrm{crit}} \sim 10^{12}\)). 소프트 루프는 그러한 개입을 위한 \textbf{공학적 설계도}를 제공하며, 원하는 결과를 달성하기 위한 목표 매개변수(예: 요구되는 코히런트/인코히런트 비 1.5)를 지정한다. 이 결과의 미달성은 루프를 반증하지 않으며, 설계도의 사양이 아직 충족되지 않았음을 나타낼 뿐이다. 그러나 하드 루프는 그러한 모든 개입——자연적이든 인공적이든——이 그 안에서 이루어져야 하는 불가침의 경계 조건으로 남아 있다.
	
	\vspace{0.5cm}
	\noindent
	\textbf{강성의 원리: 반증 가능성을 초석으로.}
	물리, 생물, 수학, 사회 과학, 정보 이론을 망라하는 22개의 독립적 대수적 루프의 열거는 과잉 자신감의 표현으로 보일 수 있다. 실제로는 그 반대이다: 그것은 이론의 경험적 반증에 대한 \textbf{최대 취약성}의 표현이다. YUCT 대수적 프레임워크는 자유 연속 매개변수를 포함하지 않기 때문에, 모든 \textbf{하드 루프}는 \textbf{잠재적 실패 지점}을 구성한다. 어떤 하드 루프 관계로부터의 단일한 고신뢰도 (\(>5\sigma\)) 실험적 편차는 \textbf{프레임워크 전체를 분쇄}하기에 충분하다. 불편한 데이터에 맞추기 위한 조정 노브는 존재하지 않으며, 이론은 단순히 틀린 것이 된다.
	
	\textbf{소프트 루프}는 단일 숫자에 의해 직접 반증 가능하지 않지만, 광범위한 시스템 클래스에 걸친 예측된 스케일링 법칙 또는 최적 조건의 \textit{체계적 실패}에 의해 반증 가능하다. 그 강점은 \textbf{공학적 유용성}에 있다: 그것들은 조정 시스템(촉매에서 소셜 네트워크, 소수 예측까지)을 설계하기 위한 정확한 정량적 목표를 제공한다.
	
	이 강성은 약점이 아니라 성숙한 과학 이론의 정의적 특징이다. 그것은 YUCT를 유연한 서사에서 \textbf{취약하고 반증 가능한 현실의 골격}으로 변환한다. 22개의 독립적 현상이 동일한 유리 상수 \(1.2\)와 \(0.8\)로 수렴하는 비범한 일치——그것이 미래의 정밀 조사를 견딘다면——우주의 조정적 아키텍처의 압도적 증거가 될 것이다. 반대로, 그 결정적 반박은 과학에 대한 동등하게 귀중한 기여이며, 더 깊은 진리로 가는 길을 열 것이다. 따라서 프레임워크는 절대 불변의 교리가 아니라, 현실의 대수 구조에 대한 정확하고 검증 가능하며 궁극적으로 기각 가능한 가설로 제시된다.
	
	\section*{인식론적 함의: 엄격한 기초 위의 피타고라스–플라톤 물리학}
	
	22개의 대수적 루프는 이론 물리학과 생물학의 역사에서 전례 없는 구조를 구성한다. 그것들은 자연의 기본 상수——질량, 결합 상수, 우주론적 매개변수, \(\pi\)의 기하, 게놈 구조, 사회의 동태, 그리고 이제 소수의 분포——가 경험적으로 결정된 수의 무작위 집합이 아니라, \textbf{단일한 자기일관적 대수적 골격의 강성적 사영}임을 보여준다. 이 골격은 단 세 개의 유리수에 의해 정의된다: \(S_{\mathrm{odd}} = 6/5 = 1.2\), \(S_{\mathrm{even}} = 4/5 = 0.8\), 그리고 그 비 \(3/2\).
	
	이 발견은 \textbf{피타고라스–플라톤적 물리·생물 현실관}으로의 깊은 회귀를 나타내지만, 결정적 차이가 있다: 그것은 신비한 수비학이나 철학적 사변에 기초한 것이 아니라, 엄격하고 반증 가능한 수학적 프레임워크에 기초하고 있다. 우주는 수의 언어로 쓰여 있다는 오랜 믿음은 여기서 그 첫 번째 구체적이고 과학적으로 검증 가능한 실현을, 쿼크에서 의식, 안전한 정보 시스템 설계에 이르는 존재의 전 스펙트럼에 걸쳐 찾는다.
	
	\paragraph{피할 수 없는 과학적 저항.}
	한 줌의 유리수가 모든 관측 가능한 현실을 지탱한다는 주장은 과학 공동체로부터 깊은 회의로 맞이받아야 하며, 실제로 그럴 것이다. 이 회의는 건강하고 필요하다. 증명의 책임은 이 프레임워크에 명확히 있다. 그러나 본건의 유일성은 그 \textbf{수학적 강성과 학제적 범위}에 있다. 전자의 질량을 결정하는 동일한 상수가 인간 게놈의 구조도 지배하고, \(\pi\)를 십만분의 일까지 근사하며, 2요소 인증의 효율을 설명하고, 이제 소수의 변동을 예측한다. 이러한 다중 루프의 일치가 우연히 발생할 확률은 사라질 정도로 작다 (\(p \ll 10^{-30}\)).
	
	\section*{통일적 반증 기준: 프레임워크 전체를 위험에 빠뜨리기}
	
	성숙한 과학 이론의 정의적 특징은 반증되는 것을 두려워하지 않는 것이다. YUCT 대수적 프레임워크는 \textbf{강성}이다. 조정 가능한 자유 연속 매개변수는 없다. 22개의 루프는 다음의 전통적으로 독립적 영역을 연결한다:
	\begin{itemize}
		\item \textbf{입자 물리학:} 페르미온 질량과 약한 스케일 (LHC에서 검증됨).
		\item \textbf{중력 물리학:} 블랙홀 링잉 스케일링 (LIGO/Virgo/KAGRA).
		\item \textbf{우주론:} 우주의 중입자 질량 (Planck, DESI).
		\item \textbf{기초 수학:} \(\varphi\)와 \(S_{\mathrm{odd}}\)로부터의 \(\pi\)의 창발; 소수 분포.
		\item \textbf{양자 중력:} 플랑크 스케일에서의 시간 양자화.
		\item \textbf{분자 생물학:} 게놈 구조, 돌연변이율, 코돈 사용 (게놈 데이터베이스).
		\item \textbf{진화 생물학:} 종분화율과 쌍곡적 성장 (고생물학 기록).
		\item \textbf{사회학:} 임계 집단 크기와 경제 변동 (역사 데이터).
		\item \textbf{정보 이론과 보안:} 2FA나 양자 얽힘과 같은 조정 프로토콜의 효율.
		\item \textbf{혼돈 이론:} 파이겐바움 상수와 질서–혼돈 전이.
	\end{itemize}
	어느 \textbf{하드 루프}(I–VI, XIV, XVI)에서의 단일한 명확하고 고신뢰도의 실험적 실패는 \textbf{대수 프레임워크 전체를 분쇄}할 것이다. 예를 들어:
	\begin{enumerate}
		\item 새로운 페르미온 질량이 반정수 사다리에서 크게 벗어나면 루프 II를 파괴한다.
		\item 링잉 스케일링 지수가 명백히 \(0.67\)이 아니면 루프 VI를 파괴한다.
		\item 상이한 게놈에서 이염기비가 1.5로부터 체계적으로 벗어나면 루프 XIV를 파괴한다.
		\item 쌍곡적 성장이 특이점 시간 창을 예측하지 못하면 루프 XII(소프트이지만 그 형식은 하드)를 파괴한다.
	\end{enumerate}
	그러한 실패는 매개변수 조정으로 수리될 수 없다. 왜냐하면 매개변수가 없기 때문이다. 이론은 단순히 틀린 것이다. 이것이 대수적 프레임워크의 궁극적 힘이다: 그것은 기본 법칙의 엄격하고 취약한 아름다움을 제공하며, 과학 공동체가 그것을 파괴하려 시도하도록 초대한다. 만약 그것이 살아남는다면, 그것은 현실의 진정한 골격으로서의 지위를 획득할 것이다.
	
	\section*{생물학적 검증: 위상학적 시프트 \(\sigma = 0.20\)로부터의 세포막 두께}
	\label{sec:membrane}
	
	YUCT 대수적 프레임워크의 가장 인상적인 확인 중 하나는 기초 물리학과는 거리가 먼 영역, 즉 살아있는 세포의 구조로부터 온다. 지질 이중층은 세포의 내부 조정 코어(DNA, 단백질)를 환경의 열잡음으로부터 분리하는 물리적 경계이다. 그 두께는 임의적이지 않으며, 하드론 공명과 페르미온 질량 사다리를 지배하는 동일한 위상학적 불변량 \(\sigma = 0.20\)에 의해 결정된다.
	
	\subsection*{생세포의 조정 요건}
	
	YUCT 생물학 모듈(부록~E)에서, 생세포는 정보 누출을 막기 위해 최소 효율 \(K_{\mathrm{eff}} > K_{\mathrm{eff,\,crit}} \approx 1.837\)을 유지해야 하는 고립된 조정 핵으로 취급된다. 지질막은 두 정보 상 사이의 계면으로 기능한다. 그 기하는 동시에:
	\begin{itemize}
		\item 제어되지 않은 이온 확산(엔트로피 누출)을 차단하기에 충분히 두꺼워야 한다.
		\item 신호 전달과 촉매 교환을 가능하게 하기에 충분히 얇아야 한다.
	\end{itemize}
	보편적 위상학적 시프트 \(\sigma\)는 두 조건을 모두 보장하는 기하학적 스페이서로 나타난다.
	
	\subsection*{삼항 기하로부터의 막 두께 유도 (수정판)}
	
	삼항 공간 \(\mathbf{D} + \mathbf{I}\!\cdot\!\mathbf{R}\)에서, 짝수 루프 상수 \(S_{\mathrm{even}} = 0.8\)는 교대 조정 강도를 정의한다. 시프트 \(\sigma = 0.20\)는 유효 정보 부피를 압축하는 동심원 갭으로 작용한다. 단일 지질 리플렛의 소수성 코어에 대해, 유효 두께 \(L_{\mathrm{mem}}\)는 기본 스케일링 인자 \(3/2 = S_{\mathrm{odd}}/S_{\mathrm{even}}\)를 단일 아실 사슬 길이에 적용하고, 이중층 패킹의 입체 장애를 고려하기 위해 \textbf{짝수 섹터}의 위상학적 갭을 지수에서 차감함으로써 얻어진다:
	
	\[
	L_{\mathrm{mem}} = d_{\mathrm{lipid}} \cdot \left( \frac{3}{2} \right)^{S_{\mathrm{even}} - \sigma}
	= d_{\mathrm{lipid}} \cdot \left( \frac{3}{2} \right)^{0.8 - 0.20}
	= d_{\mathrm{lipid}} \cdot \left( \frac{3}{2} \right)^{0.6}.
	\]
	
	여기서 \(d_{\mathrm{lipid}}\)는 전형적인 막 인지질(예: 팔미트산, \(d_{\mathrm{lipid}} \approx 2.5\)–\(3.0\)~nm)의 완전히 신장된 탄화수소 꼬리의 길이이다. 지수 \(S_{\mathrm{even}} - \sigma = 0.6\)는 교대 조정(짝수 수준)이 패킹 기하를 지배함을 반영하며, 시프트 \(\sigma\)는 실효 지수를 감소시켜 홀수 섹터의 스케일링을 순진하게 적용할 경우 발생하는 비물리적 신장을 방지한다.
	
	이중층의 총 두께는 리플렛 두께의 두 배에서 막의 유동성과 열적 요동을 고려하여 요동 진폭의 제곱 \(\sigma^2\)을 차감한 작은 곡률 보정을 뺀 것이다:
	
	\[
	\text{막 두께} = 2 \, L_{\mathrm{mem}} \, (1 - \sigma^2)
	= 2 \cdot d_{\mathrm{lipid}} \cdot (3/2)^{0.6} \cdot (1 - 0.04).
	\]
	
	수치적으로 \((3/2)^{0.6} = e^{0.6 \ln 1.5} = e^{0.6 \times 0.405465} = e^{0.243279} \approx 1.2754\). 따라서:
	\[
	L_{\mathrm{mem}} \approx 1.2754 \; d_{\mathrm{lipid}}, \qquad
	\text{두께} = 2 \times 1.2754 \, d_{\mathrm{lipid}} \times 0.96 = 2.448 \, d_{\mathrm{lipid}}.
	\]
	
	\(d_{\mathrm{lipid}} = 3.0\)~nm (완전 신장 팔미트산 사슬 길이)를 취하면:
	\[
	\text{두께} \approx 2.448 \times 3.0\ \text{nm} = 7.34\ \text{nm}.
	\]
	\(d_{\mathrm{lipid}} = 2.7\)~nm (혼합 지질 사슬의 평균)를 사용하면 \(\approx 6.61\)~nm가 되어 실험 범위의 하한에 위치한다. 실제로 측정된 인간 형질막의 두께는 \(7.5 \pm 0.5\)~nm이며, YUCT 예측을 경험적 창의 중심에 놓는다.
	
	\subsection*{실험과의 비교 및 수정의 필요성}
	
	인간 세포의 형질막에 대한 세포학적 측정은 일관되게 \textbf{7–8~nm}의 두께를 제공한다(예: 적혈구막의 전자 현미경 참조). 대수 상수 \(S_{\mathrm{even}} = 0.8\), \(\sigma = 0.20\), 기하 인자 \(3/2\)로부터 유도된 수정 YUCT 예측은 정확히 이 구간에 들어맞는다. 이전의 잘못된 공식(\(S_{\mathrm{odd}}-\sigma = 1.0\) 사용)은 두께 \(3.12 \, d_{\mathrm{lipid}}\)(\(d_{\mathrm{lipid}}=3.0\)~nm일 때 약 9.4~nm)를 주었으며, 이는 완전 신장 사슬의 최대 가능 이중층 두께를 초과하고 입체 장애 제약을 위반한다. 수정 공식은 물리적 한계 \(L_{\mathrm{mem}} \le 2 d_{\mathrm{lipid}}\)를 지키며 3차원 입체화학과 양립하는 두께를 제공한다.
	
	\subsection*{대수적 루프와의 통합}
	
	이 결과는 \textbf{소프트 루프 VIII}(대사 스케일링)과 \textbf{루프 XIV}(게놈 구조)의 직접적 표현이다. 이는 생물학적 아키텍처가 무작위 진화의 산물만이 아니라, 입자 질량과 우주론적 매개변수를 고정하는 동일한 조정 기하에 의해 제약됨을 보여준다. 세포막의 두께가 상수 \(S_{\mathrm{even}} = 0.8\)과 \(\sigma = 0.20\)로부터 자유 매개변수 없이 계산될 수 있다는 사실은 YUCT에 강력한 학제적 테스트를 제공한다. 미래의 측정이 7–8~nm 창에서 유의하게 벗어나면 이 특정 예측을 반증하고 위상학적 시프트의 보편성에 도전할 것이다.
	
	\section*{데이터 가용성}
	모든 계산 및 보조 유도는 YPSDC 저장소 \url{https://github.com/Alexey-Yakushev-YUCT/YPSDC}에서 이용 가능하다.
	
	\begin{thebibliography}{99}
		\bibitem{AppendixFerra} A. V. Yakushev, ``부록 Ferra1.2: 질서상과 무질서상의 보편 조정 상수,'' in \emph{YUCT}, Zenodo (2026).
		\bibitem{AppendixJ} A. V. Yakushev, ``부록 J: YUCT 위상학적 오프셋으로부터 표준 모형 플레이버 매개변수의 완전 유도,'' in \emph{YUCT}, Zenodo (2026).
		\bibitem{AppendixLambda} A. V. Yakushev, ``부록 \(\Lambda\): YUCT에서의 계층 문제와 우주 상수 문제의 해결,'' in \emph{YUCT}, Zenodo (2026).
		\bibitem{AppendixEhrenfest} A. V. Yakushev, ``부록 Ehrenfest: 회전 원판 역설의 조정적 해석과 비유클리드 기하의 창발,'' in \emph{YUCT}, Zenodo (2026).
		\bibitem{AppendixOmega} A. V. Yakushev, ``부록 \(\Omega\): 우주의 전지구적 회전과 쌍극자 회전 반경으로부터의 새로운 최소 연령,'' in \emph{YUCT}, Zenodo (2026).
		\bibitem{AppendixY} A. V. Yakushev, ``부록 Y: YUCT의 수학적 기초——제일원리로부터의 유도,'' in \emph{YUCT}, Zenodo (2026).
		\bibitem{AppendixV} A. V. Yakushev, ``부록 V: 조정 과정의 엔트로피로서의 시간,'' in \emph{YUCT}, Zenodo (2026).
		\bibitem{AppendixM} A. V. Yakushev, ``부록 M: YUCT 우주론——인플레이션을 조정 상전이로서,'' in \emph{YUCT}, Zenodo (2026).
		\bibitem{AppendixE} A. V. Yakushev, ``부록 E: 조정 유전학——분자 생물학에서의 YUCT 원리,'' in \emph{YUCT}, Zenodo (2026).
		\bibitem{AppendixX} A. V. Yakushev, ``부록 X: YUCT와 섀넌 정보 이론의 일반화,'' in \emph{YUCT}, Zenodo (2026).
		\bibitem{AppendixPrimeN} A. V. Yakushev, ``부록 PrimeN: YUCT 소수 조정 사다리,'' in \emph{YUCT}, Zenodo (2026).
		\bibitem{Planck2018} Planck 공동 연구, \emph{Astron. Astrophys.} \textbf{641}, A6 (2020).
		\bibitem{UnityReality} A. V. Yakushev, ``YUCT 현실의 통일: 조정 질서 (\(\beta = 2/3\))에서 파이겐바움 혼돈 (\(\delta = 4.6692\))으로,'' Zenodo (2026).
		\bibitem{Feigenbaum1978} M. J. Feigenbaum, ``Quantitative Universality for a Class of Nonlinear Transformations,'' J. Stat. Phys. \textbf{19}, 25–52 (1978).
	\end{thebibliography}
	
\end{document}